\chapter{Evaluation: Joint Multipath Optimization}
\label{ch:p2eval}

This chapter evaluates the Part~II system. \Cref{sec:p2eval:setup} gives the setup and baselines;
\cref{sec:p2eval:results} reports the central finding -- that the value of learned multipath scheduling
depends sharply on how volatile the network is; and \cref{sec:p2eval:ablation} isolates the contributions
of the joint design and the scoring transport agent. The numbers reported here are from preliminary runs
on the mock backend; packet-level NS-3 confirmation is pending (\cref{sec:p2impl:gaps}) and is marked
where relevant.

\section{Setup and Baselines}
\label{sec:p2eval:setup}

\paragraph{Scenarios.} We use two contrasting scenarios. The \emph{static} scenario is the non-dominant
four-path topology of \cref{sec:p2impl:rt}, with stationary (if heavily loaded) paths. The \emph{dynamic}
scenario is a six-path setting with the non-stationary dynamics of \cref{sec:p2impl:dataplane} enabled --
path churn, regime shifts, and congestion bursts -- so that the identity and quality of the best paths
change during a call.

\paragraph{Metric.} The primary metric is mean per-window QoE (\cref{eq:p2qoe}); we also report the
underlying achieved bitrate, latency, loss, and VMAF.

\paragraph{Baselines.} The learned pair (application agent + transport agent) is compared against fixed
scheduling policies driving the \emph{same} learned bitrate controller where applicable: \emph{even}
split across live paths, \emph{single} best-active path, and \emph{proportional} allocation by capacity.
The decisive ablation is \emph{app-only}: the learned bitrate controller with a naive even split,
isolating the value of the learned transport agent.

\section{Results: When Does Splitting Help?}
\label{sec:p2eval:results}

The headline result is that the learned split's value is conditional on volatility.
\Cref{tab:p2results} summarizes preliminary QoE across scenarios and methods.

\begin{table}[t]
  \centering
  \caption{Preliminary mean QoE (\cref{eq:p2qoe}) by scenario and method (mock backend). Higher is
  better. In the static scenario the split barely matters -- the learned transport agent adds little over
  app-only -- whereas in the dynamic scenario the split is decisive and app-only collapses.}
  \label{tab:p2results}
  \begin{tabular}{lcc}
    \toprule
    Method & Static 4-path & Dynamic 6-path \\
    \midrule
    Learned pair (app + transport) & \textbf{0.633} & \textbf{0.606} \\
    App-only (learned rate, even split) & 0.665 & $-0.212$ \\
    Single best-active path & 0.514 & 0.402 \\
    Proportional (by capacity) & \gap{fill} & 0.209 \\
    Even split & \gap{fill} & 0.109 \\
    \bottomrule
  \end{tabular}
\end{table}

Two observations stand out. First, in the \textbf{static} scenario the learned pair (0.633) beats the
single-best-path baseline (0.514) by roughly 23\%, but the app-only ablation (0.665) is actually on par
with -- even marginally above -- the full system. When paths are stable, a good bitrate controller over a
simple split captures nearly all the achievable QoE, and the transport agent has little left to earn. A
practitioner reading only the static numbers might reasonably conclude that learned scheduling is not
worth its complexity.

Second, the \textbf{dynamic} scenario overturns that conclusion. The learned pair (0.606) again leads,
but now the app-only ablation \emph{collapses to $-0.212$}: with paths churning and bursting, a fixed even
split routes traffic onto dead or congested paths, driving latency and loss up despite a well-chosen
bitrate. The fixed-schedule baselines degrade in order of how much they react to path state -- single-best
(0.402), proportional (0.209), even (0.109) -- and only the learned transport agent, which observes path
liveness and reallocates, sustains high QoE. This is the condition under which learned multipath
scheduling is not a marginal optimization but the difference between a usable and an unusable call.

\gap{add variance/confidence intervals over seeds, and a time series from a representative dynamic
episode showing the transport agent reallocating away from a path as it churns or bursts (the
\texttt{dyn-scoring} run illustrates convergence from strongly negative early-episode reward to stable
positive QoE with latency in the tens of milliseconds and loss a few percent).}

\placeholderfig{Figure: QoE over training episodes for the learned pair vs.\ baselines in the dynamic
scenario, and a within-episode trace of the per-path split tracking path liveness.}

\section{Ablation Study}
\label{sec:p2eval:ablation}

\paragraph{Joint vs.\ app-only.} \Cref{tab:p2results} already provides the core ablation: removing the
learned transport agent (app-only) is nearly free when the network is stable but catastrophic when it is
volatile. This quantifies the design premise of \cref{sec:p2design:motivation} -- the two decisions must
be co-optimized precisely in the regimes where the network is hard.

\paragraph{Flat vs.\ scoring transport.} We compare the flat transport agent (fixed path count) against
the permutation-equivariant scoring agent (\cref{sec:p2design:scoring}) on the dynamic scenario, where the
number of live paths changes.
\gap{report the flat-vs-scoring comparison; the expected finding is that only the scoring/mask-aware agent
handles a changing path count gracefully, since the flat agent cannot represent appearing/disappearing
paths.}

\paragraph{Backend parity.} All numbers above are from the mock backend. Confirming them on the
packet-level NS-3 backend requires the dynamics port of \cref{sec:p2impl:gaps}.
\gap{report NS-3 results for the static scenario now, and for the dynamic scenario once the C++ dynamics
land, to establish that the mock findings transfer to packet-level simulation.}

\section{Summary}
\label{sec:p2eval:summary}
Preliminary results support the joint, multipath design: co-optimizing bitrate and per-path split matches
or beats every fixed-scheduling baseline, and the advantage of learned scheduling grows from negligible
in a stable network to decisive under churn and bursts. The main open items are seed-level statistics,
the flat-vs-scoring comparison, and NS-3 confirmation.
